\documentclass[10pt,a4paper]{amsart}
\usepackage[margin=19mm]{geometry}
\usepackage{amsmath,amssymb,mathtools}
\usepackage[scr=rsfs,cal=euler]{mathalfa}
\usepackage{xfrac}
\usepackage[unicode,hidelinks,bookmarks=false]{hyperref}
\newcommand{\ie}{i.e.\ }
\newcommand{\cf}{cf.\ }
\newcommand{\resp}{respectively\ }
\newcommand{\Subsection}{\subsection}
\providecommand{\nobreakdash}{\nobreak\hbox{-}}
\setlength{\emergencystretch}{2em}
\multlinegap0pt
\usepackage{bm}
\usepackage{amssymb}
\usepackage{algorithm}\usepackage{algpseudocode}
\newtheorem{lemma}{Lemma}
\title{LABS: Extending the scope of binary segmentation via a look-ahead device}
\author{Piotr Fryzlewicz}
\date{Source: arXiv:2608.21122v1 (CC BY 4.0)}
\begin{document}
\maketitle

\noindent\emph{Source and licence.} LABS: Extending the scope of binary segmentation via a look-ahead device. Version arXiv:2608.21122v1 (CC BY 4.0), distributed by arXiv under the Creative Commons Attribution 4.0 International licence, http://creativecommons.org/licenses/by/4.0/, as recorded in the arXiv licence metadata for this exact version and shown on the abstract page https://arxiv.org/abs/2608.21122v1. Full licence notice: \texttt{licenses/arxiv-2608.21122-CC-BY-4.0.txt}.\par\smallskip

\noindent\textit{Editorial scope.} Counted unit: the article's lemma lem:loc (source0.tex lines 1711--1736). The article's complete original proof of it is delivered with it (source0.tex lines 1738--1830, one \texttt{proof} environment of 93 lines). No mathematics is written, repaired or re-derived: the statement and the proof are the article's own lines, in the article's own notation, and no retained line is substituted or re-flowed.  Retained uncounted as the source-local context the proof needs: lines 850--872; lines 967--979; lines 1013--1023; lines 1056--1105.  Nothing was omitted from any retained span, so the package carries no omission record.  No retained line contains a citation, so the wrapper carries no bibliography and no bibliography entry.  No retained line contains a figure environment, an image-inclusion command or drawing code, so no image is delivered and the package declares no asset dependency.  Editorial formatting only, in addition: an AMS \texttt{amsart} preamble that restates the article's own declarations and macros used by the retained lines, namely the environments \texttt{lemma} on separate counters, together with the standard packages those lines need.  The retained environments print in this document as Lemma 1 in order of appearance, whereas the article prints the same items with the numbering of its own full document.  The complete native source of the article as distributed in the arXiv e-print for this exact version is bundled unchanged as \texttt{sources/208227/source0.tex}.
\begin{lemma}[Representation]
\label{lem:rep}
Let $V(s,e) := \{ j : s+2 \le \tau_j \le e-1 \}$. Then:
\begin{enumerate}
\item[\emph{(a)}] there are $a,c \in \mathbb{R}$ with
$f_t = a + ct + \sum_{j \in V(s,e)} \Delta_j (t-\tau_j)_+$ for $t \in W$, and
hence
$\langle f, \hat\psi_b\rangle = \sum_{j \in V(s,e)} \Delta_j \langle
\psi_{\tau_j}, \hat\psi_b\rangle$ for every $b \in B_{s,e}$;
\item[\emph{(b)}] $\max_{b} |\langle f, \hat\psi_b\rangle| \le \sum_j
\ell_j(s,e)$, and for any $U$ and any $b, b' \in B_{s,e}$,
\[
\Big| \sum_{j \notin U} \Delta_j \langle \psi_{\tau_j},
\hat\psi_b - \hat\psi_{b'} \rangle \Big|
\;\le\; \Big( \sum_{j \notin U} \ell_j(s,e) \Big)\,
\|\hat\psi_b - \hat\psi_{b'}\| ;
\]
\item[\emph{(c)}] if $V(s,e) = \emptyset$ then $\langle f,\hat\psi_b\rangle = 0$
for every $b$; and if $V(s,e) = \{j\}$ then
$\max_b |\langle f,\hat\psi_b\rangle| = \ell_j(s,e)$, attained at $b = \tau_j$
and there only.
\end{enumerate}
\end{lemma}

\begin{lemma}[Uniform noise bounds]
\label{lem:noise}
Let
\[
E_n := \Big\{ \max_{s,e,b} |\langle \varepsilon, \hat\psi_b\rangle| \le
\sigma(8\log n)^{1/2} \Big\} \cap
\Big\{ \max_{s,e,b \ne b'} \frac{|\langle \varepsilon, \hat\psi_b -
\hat\psi_{b'}\rangle|}{\|\hat\psi_b - \hat\psi_{b'}\|} \le
\sigma(10\log n)^{1/2} \Big\},
\]
the maxima being over all $0 \le s < e \le n$ and $b,b' \in B_{s,e}$. Then
$\mathbb{P}(E_n) \to 1$. Write $Z_n := \sigma(8\log n)^{1/2}$.
\end{lemma}

\begin{lemma}[Two-sided separation]
\label{lem:sep}
Write $\rho(\theta,b) = |\langle \hat\psi_\theta, \hat\psi_b\rangle|$. There
are absolute constants $c_0 > 0$ and $C_0 < \infty$ such that for every window,
every $\theta$ with $x := \min(\theta-s, e+1-\theta) \ge 2$, and every
$b \in B_{s,e}$,
\[
c_0 \min\Big\{ \Big(\frac{b-\theta}{x}\Big)^2, 1 \Big\} \;\le\;
1 - \rho(\theta,b) \;\le\; C_0 \Big(\frac{b-\theta}{x}\Big)^2 .
\]
\end{lemma}

\begin{lemma}[Gram determinants]
\label{lem:gramdet}
\emph{(a) Continuum.} For $0 < u \le v < 1$, with $\Psi$ formed on $[0,1]$,
\[
\|\Psi_u\|^2\|\Psi_v\|^2 - \langle \Psi_u,\Psi_v\rangle^2
\;=\; \frac{u^3 (1-v)^3 (v-u)^2 \, Q(u,v)}{36},
\qquad Q(u,v) := 4v - u - 3uv ,
\]
and consequently
\begin{equation}
\label{eq:rhoexact}
1 - \rho(u,v)^2 \;=\; \frac{(v-u)^2\,Q(u,v)}{4\,(1-u)^3 v^3},
\qquad\text{with}\qquad
3u(1-u) \;\le\; Q(u,v) \;\le\; 4 .
\end{equation}

\emph{(b) Discrete.} Let $1 \le k < l \le m-1$ and put $a = k$, $c = l-k$ and
$b = m-l+1$, so that $a+b+c = m+1$. Then
\[
\|\psi_k\|^2\|\psi_l\|^2 - \langle\psi_k,\psi_l\rangle^2
\;=\; \frac{a\,b\,c\,(a-1)(b-1)\,P(a,b,c)}
{36\,(a+b+c)(a+b+c-1)^2(a+b+c-2)} ,
\]
where
\begin{align*}
P(a,b,c) &= 3a^2b^2c + 4a^2bc^2 - 3a^2bc + 2a^2b - 2a^2c^2 + 3a^2c - a^2 \\
&\quad + 4ab^2c^2 - 3ab^2c + 2ab^2 + 4abc^3 - 8abc^2 + 11abc - 4ab \\
&\quad - 2ac^3 + 6ac^2 - 7ac + 3a - 2b^2c^2 + 3b^2c - b^2 \\
&\quad - 2bc^3 + 6bc^2 - 7bc + 3b + c^3 - 4c^2 + 5c - 2 .
\end{align*}
Its terms of top degree are $abc(3ab+4ac+4bc+4c^2)$, and with $u = a/m$, $v = (a+c)/m$ and $m = a+b+c-1$,
\[
m^2 Q(u,v) \;=\; 3ab + 4ac + 4bc + 4c^2 - 3a - 4c ,
\]
The determinant here is the left-hand side of part (b), namely the determinant
of the $2\times2$ Gram matrix of $\psi_k$ and $\psi_l$. Each squared norm and
each inner product is of order $m^3$, so every term in that determinant is of
order $m^6$. Equivalently, the numerator in part (b) has total degree ten and
the denominator has degree four. After division by $m^6$, its highest-degree
terms give the continuum determinant in part (a), with
$u=a/m$ and $v=(a+c)/m$.

\emph{(c) Positivity.} For $1 \le b_1 \le b_2 \le m-1$, put $a = b_1$, $c = b_2-b_1$ and $d = m-b_2+1$. Then
\[
\langle \psi_{b_1}, \psi_{b_2}\rangle \;=\;
\frac{a\,d\,(a-1)(d-1)\big\{ 2ad-a-d+2 + 3c\,(a+d+c-1) \big\}}
{6\,(a+d+c-2)(a+d+c-1)(a+d+c)} ,
\]
which is strictly positive whenever $a, d \ge 2$, that is whenever $b_1, b_2 \in B_{s,e}$. In particular $\rho(b_1,b_2) = \langle\hat\psi_{b_1},\hat\psi_{b_2}\rangle$, with no absolute value needed.
\end{lemma}

\begin{lemma}[Detection and localization]
\label{lem:loc}
Let $(s,e]$ be a window and put
\[
D_U(w):=\sum_{j\in U}\Delta_j\langle\psi_{\tau_j},w\rangle,
\qquad D_U(b):=D_U(\hat\psi_b),
\qquad
\varrho := \sum_{j\notin U}\ell_j(s,e) ,
\]
for vectors $w\in\mathbb R^W$ and some
$U\subseteq\{1,\ldots,N\}$. Suppose there are $b^* \in B_{s,e}$ and constants $a_0, \varkappa', \varrho_1 > 0$ such that
\begin{enumerate}
\item[(i)] $\min(b^*-s,\, e-b^*) \ge \varrho_1 n$;
\item[(ii)] $A := |D_U(b^*)| \ge a_0 n^{1/2}$, and
$|D_U(b^*)| - |D_U(b)| \ge \varkappa' A \{(b-b^*)/n\}^2$
for every $b \in B_{s,e}$ with $|b-b^*| \ge r_0$, where $r_0 \le C_r n^{1/2}$;
\item[(iii)] $\varrho \le \sigma(\log n)^{1/2}$.
\end{enumerate}
Then there are $n_0$ and $C_1$, depending only on $a_0,\varkappa',\varrho_1,C_r,\sigma,c_0,C_0$, such that for $n \ge n_0$ on
$E_n$,
\[
\max_b \mathcal{C}_{s,e}(b) \ge \tfrac12 A > \lambda,
\qquad
|\hat b - b^*| \le C_1 \bar r_n .
\]
\end{lemma}

\begin{proof}
Define the linear functional
\[
R(w):=\sum_{j\notin U}\Delta_j\langle\psi_{\tau_j},w\rangle.
\]
Then $\langle f,w\rangle=D_U(w)+R(w)$,
$|R(w)|\le\varrho\|w\|$, and
$|R(\hat\psi_b)-R(\hat\psi_{b'})|
\le\varrho\|\hat\psi_b-\hat\psi_{b'}\|$ by Lemma~\ref{lem:rep}(b). Since $A \ge a_0n^{1/2}$ and
$\varrho + Z_n = O\{(\log n)^{1/2}\}$ we get
$\mathcal{C}(b^*) \ge A - \varrho - Z_n \ge \tfrac12 A > \lambda$, the first
claim.

If $|\hat b - b^*| < r_0$ the second claim holds, since $r_0 \le C_r n^{1/2} \le C_r \bar r_n$; assume therefore $|\hat b - b^*| \ge r_0$, so that hypothesis (ii) may be used at $b = \hat b$. Put $\varsigma^* = \mathrm{sign}\, D_U(b^*)$, $u^* = \varsigma^*\hat\psi_{b^*}$,
and $u = \varsigma \hat\psi_{\hat b}$ with $\varsigma$ chosen so that
$\langle X,u\rangle = \mathcal{C}(\hat b) \ge 0$. From
$\langle X,u\rangle \ge \langle X,u^*\rangle$ and Lemma~\ref{lem:noise},
\begin{equation}
\label{eq:basic}
\langle f,u^*\rangle - \langle f,u\rangle \;\le\;
\langle \varepsilon, u - u^*\rangle .
\end{equation}
Since $D_U(u^*)=|D_U(b^*)|$ by the choice of $\varsigma^*$, and
$D_U(u)\le|D_U(\hat b)|$,
\[
\langle f,u^*\rangle-\langle f,u\rangle
\ge |D_U(b^*)|-|D_U(\hat b)|-|R(u^*-u)|
\ge |D_U(b^*)|-|D_U(\hat b)|-\varrho\|u-u^*\|.
\]
Combining this inequality with \eqref{eq:basic} and hypothesis (ii), and
writing $\Delta:=|\hat b-b^*|$, gives
\begin{equation}
\label{eq:gapineq}
\varkappa' A(\Delta/n)^2
\le |D_U(b^*)|-|D_U(\hat b)|
\le \varrho\|u-u^*\|+\langle\varepsilon,u-u^*\rangle.
\end{equation}

The noise term is treated differently before and after Step~2 proves that
$\varsigma=\varsigma^*$. The increment bound is the second event in the
definition of $E_n$; it controls
$\langle\varepsilon,\hat\psi_b-\hat\psi_{b'}\rangle$, a difference of
unsigned contrast vectors. If $\varsigma=-\varsigma^*$, then
$u-u^*=\pm(\hat\psi_{\hat b}+\hat\psi_{b^*})$ is instead a sum, so the second
event does not apply. Before Step~2 excludes this case, we use the pointwise
part of $E_n$, namely its first event, separately at $\hat b$ and $b^*$:
\[
\big| \langle\varepsilon, u-u^*\rangle \big| \;\le\; \big|\langle\varepsilon,u\rangle\big| + \big|\langle\varepsilon,u^*\rangle\big| \;\le\; 2Z_n ,
\]
This inequality holds for either relation between the signs. After Step~2 has
proved $\varsigma=\varsigma^*$,
$u-u^*=\pm(\hat\psi_{\hat b}-\hat\psi_{b^*})$, and the second event in $E_n$
can be used. Lemma~\ref{lem:sep} then makes both terms on the right of
\eqref{eq:gapineq} proportional to $\Delta/n$, whereas its left-hand side is
proportional to $(\Delta/n)^2$. For $\Delta>0$, one factor $\Delta/n$ can then
be canceled. The bound $|R(w)|\le\varrho\|w\|$ holds for every $w$ and does
not depend on the relation between the signs.

\emph{Step 1 (preliminary bound).} Using $\|u-u^*\| \le 2$, the pointwise bound just described, and (iii), $\varkappa' A(\Delta/n)^2 \le 2\varrho + 2Z_n \le 2\sigma(\log n)^{1/2} + 2\sigma(8\log n)^{1/2} \le 8\sigma(\log n)^{1/2}$, so $\Delta/n = O\{n^{-1/4}(\log n)^{1/4}\} \to 0$. The second, increment event in $E_n$ is not used at this stage.

\emph{Step 2 (sign of the inner product).} By Step~1 and Lemma~\ref{lem:sep} with $x = \min(b^*-s,e-b^*) \ge \varrho_1 n$ we get $\rho := \rho(b^*,\hat b) \to 1$. Suppose $\langle u,u^*\rangle = -\rho$. Then $\|u+u^*\|^2 = 2-2\rho \to 0$, while $\|\Pi^\perp_W f\| \le \sum_j \ell_j \le N\bar d\, n^{1/2}/\sqrt2$, so
$|\langle f,u\rangle + \langle f,u^*\rangle| \le Cn^{1/2}\sqrt{2-2\rho} =
o(A)$. Since $\langle f,u^*\rangle \ge A - \varrho \ge \tfrac34 A$, this forces
$\langle f,u\rangle \le -\tfrac12 A < 0$; but
$\langle f,u\rangle \ge \mathcal{C}(\hat b) - Z_n \ge \tfrac12 A - Z_n > 0$, a
contradiction. Hence $\langle u,u^*\rangle = \rho > 0$. This does not yet give $\varsigma = \varsigma^*$: a positive signed correlation could equally arise from opposite multipliers applied to negatively correlated vectors. What excludes that is Lemma~\ref{lem:gramdet}(c), by which $\langle\hat\psi_{\hat b},\hat\psi_{b^*}\rangle > 0$ for any two admissible candidates; since $\langle u,u^*\rangle = \varsigma\varsigma^*\langle\hat\psi_{\hat b},\hat\psi_{b^*}\rangle$, it follows that $\varsigma = \varsigma^*$. Consequently $\|u-u^*\| = \{2(1-\rho)\}^{1/2} \le (2C_0)^{1/2}\varrho_1^{-1}\Delta/n$.

\emph{Step 3 (conclusion).} The signs now agree, so
$u-u^*=\pm(\hat\psi_{\hat b}-\hat\psi_{b^*})$ and the increment event in
$E_n$ gives
$\langle\varepsilon,u-u^*\rangle
\le\sigma(10\log n)^{1/2}\|u-u^*\|$. Substituting this and the bound of
Step~2 into \eqref{eq:gapineq} gives
\[
\varkappa' A\Big(\frac{\Delta}{n}\Big)^2
\le \big\{\varrho+\sigma(10\log n)^{1/2}\big\}
(2C_0)^{1/2}\varrho_1^{-1}\frac{\Delta}{n}.
\]
For $\Delta>0$, canceling one factor of $\Delta/n$ gives
\[
\frac{\Delta}{n}
\le\frac{(2C_0)^{1/2}}{\varrho_1\varkappa' A}
\big\{\varrho+\sigma(10\log n)^{1/2}\big\};
\]
the conclusion is immediate when $\Delta=0$. By (iii), the quantity in braces
is at most $5\sigma(\log n)^{1/2}$, and $A\ge a_0n^{1/2}$. Therefore
$\Delta\le C_1\bar r_n$ with
\[
C_1=\max\{C_r,\ 5(2C_0)^{1/2}\sigma/(\varrho_1\varkappa'a_0)\}.
\]
The term $C_r$ covers the case $\Delta<r_0$ considered before
\eqref{eq:basic}; the second term covers $\Delta\ge r_0$.
\end{proof}

\end{document}
